\documentclass{article}

\usepackage{amsmath,amssymb}
\usepackage{xurl}
\usepackage[hidelinks]{hyperref}

% Shared commands for notes in docs/paper-gaps/.

\DeclareUnicodeCharacter{2080}{\ifmmode{}_{0}\else\textsubscript{0}\fi}
\DeclareUnicodeCharacter{2081}{\ifmmode{}_{1}\else\textsubscript{1}\fi}
\DeclareUnicodeCharacter{2082}{\ifmmode{}_{2}\else\textsubscript{2}\fi}
\DeclareUnicodeCharacter{2099}{\ifmmode{}_{n}\else\textsubscript{n}\fi}

\newcommand{\Lean}{\textsc{Lean}}
\ExplSyntaxOn
\NewDocumentCommand{\leanid}{m}{
  \ifmmode
    \textsf{\detokenize{#1}}
  \else
    \group_begin:
    \urlstyle{sf}
    \tl_set:Nn \l_tmpa_tl {#1}
    \tl_replace_all:Nnn \l_tmpa_tl {\_} {_}
    \exp_args:NV \path \l_tmpa_tl
    \group_end:
  \fi
}
\ExplSyntaxOff
\providecommand{\tr}{\operatorname{tr}}
\newcommand{\tnleanIssueBase}{https://github.com/LionSR/TNLean/issues/}
\newcommand{\tnleanPrBase}{https://github.com/LionSR/TNLean/pull/}
\newcommand{\tnleanDocBase}{https://sirui-lu.com/TNLean/docs/}
\newcommand{\ghissue}[1]{\href{\tnleanIssueBase#1}{GitHub issue \##1}}
\newcommand{\ghpr}[1]{\href{\tnleanPrBase#1}{PR \##1}}
\newcommand{\doclink}[1]{\href{\tnleanDocBase#1.html}{\path{#1}}}

% Dirac notation and the identity operator, as in blueprint/src/macros/common.tex.
% \providecommand keeps note-local definitions authoritative.
\usepackage{dsfont}
\providecommand{\ket}[1]{|#1\rangle}
\providecommand{\bra}[1]{\langle #1|}
\providecommand{\braket}[2]{\langle #1 | #2 \rangle}
\providecommand{\ketbra}[2]{|#1\rangle\!\langle #2|}
\providecommand{\Id}{\mathds{1}}
\providecommand{\one}{\mathds{1}}
\providecommand{\C}{\mathbb{C}}
\providecommand{\MN}[1]{M_{#1}(\C)}
\providecommand{\MD}{M_D(\C)}

% External referencing for paper sources.  The build helper writes sanitized
% auxiliary files under build/paper-aux/ and excludes duplicated source labels.
\usepackage{xr}

% Import a generated paper auxiliary file with a caller-supplied label prefix.
% The candidate paths support builds from docs/paper-gaps/, the repository root,
% and build/paper-aux/ itself.
\newcommand{\importpaperaux}[2]{%
  \IfFileExists{../../build/paper-aux/#2.aux}{%
    \externaldocument[#1]{../../build/paper-aux/#2}%
  }{%
    \IfFileExists{build/paper-aux/#2.aux}{%
      \externaldocument[#1]{build/paper-aux/#2}%
    }{%
      \IfFileExists{#2.aux}{%
        \externaldocument[#1]{#2}%
      }{}%
    }%
  }%
}

% General external reference macros
\newcommand{\paperref}[2]{\ref{#1-#2}}
\newcommand{\papereqref}[2]{(\ref{#1-#2})}

% CPSV16 source (1606.00608 / MPDO-22-12-17-2.tex) external document and shortcuts
\importpaperaux{cpsv16-}{1606.00608/MPDO-22-12-17-2-tnlean}
\newcommand{\cpsvref}[1]{\paperref{cpsv16}{#1}}
\newcommand{\cpsveqref}[1]{\papereqref{cpsv16}{#1}}

% Verdict marker: \gapnote{<kind>}{<status>}. Kind is the policy
% classification (clarification, local-correction, scope-restriction,
% unfaithful, false-source, open-gap); status is open, wip (elimination
% underway), resolved, or historical. Machine-read by the site index;
% prints a verdict line.
\newcommand{\gapnote}[2]{\par\noindent\textbf{Verdict:} #1 (#2).\par}


\title{The Block Index in the Positive Generalized-Cocycle Corollary}
\author{}
\date{2026-08-29}

\begin{document}
\maketitle
\gapnote{local-correction}{resolved}

\section{Source context}

Let $G$ be a finite group acting on a set $\mathsf X$ of block labels. For a
block-dependent one-cochain $\chi\colon\mathsf X\times G\to\C^\times$ the
generalized coboundary is
\begin{align}
    (d\chi)^x_{g,h}
        &=\frac{\chi^{h x}_g\chi^x_h}{\chi^x_{gh}},
    \label{eq:fbc25_positive_block_index_coboundary}
\end{align}
and a block-dependent two-cochain $\Omega$ is a generalized cocycle when
$d\Omega=1$, that is, when
$\Omega^x_{g,hk}\Omega^x_{h,k}=\Omega^{k x}_{g,h}\Omega^x_{gh,k}$.

The finite-group determinant lemma in
Ref.~\cite[lines~5931--5947]{FrancoRubio2025MPUGauging} starts with a
block-dependent generalized cocycle $\Omega^x_{g,h}$ and constructs a
block-dependent cochain $\chi^x_g=\det X^x_g$. The corollary that follows it
instead asserts the existence of a positive $\chi_g$ with $\Omega=d\chi$,
written without the block label
\cite[Corollary~\texttt{cor:positive\_trivial}, lines~5950--5952]
{FrancoRubio2025MPUGauging}.

The omitted superscript is shorthand, not a restriction to a block-independent
cochain. The application to $\zeta$ in
Ref.~\cite[line~6018]{FrancoRubio2025MPUGauging} concerns an ordinary scalar
cochain. The later application to the block-dependent cochain $|L|$
\cite[lines~6023--6029]{FrancoRubio2025MPUGauging} suppresses its block label
while using the generalized conclusion block by block.

\section{A two-block counterexample}

Let $G=C_2=\{e,s\}$, with $s^2=e$, act trivially on the two-element block set
$\mathsf X=\{a,b\}$. Define a positive block-dependent $1$-cochain $\rho$ by
\begin{align}
    \rho^a_e=\rho^b_e&=1,
    &\rho^a_s&=1,
    &\rho^b_s&=2,
    \label{eq:fbc25_positive_block_index_cochain}
\end{align}
and set $\Omega=d\rho$. Because the action is trivial,
Eq.~\ref{eq:fbc25_positive_block_index_coboundary} reads
$\Omega^x_{g,h}=\rho^x_g\rho^x_h/\rho^x_{gh}$. Every value of $\Omega$ is
therefore a product and quotient of the positive numbers in
Eq.~\ref{eq:fbc25_positive_block_index_cochain}, hence positive; explicitly,
for $x\in\{a,b\}$ and $g\in G$,
\begin{align}
    \Omega^x_{e,e}=\Omega^x_{e,g}=\Omega^x_{g,e}
        &=\rho^x_e=1,
    \label{eq:fbc25_positive_block_index_identity_axes}\\
    \Omega^a_{s,s}
        &=\frac{\rho^a_s\rho^a_s}{\rho^a_e}=1,
    \label{eq:fbc25_positive_block_index_value_a}\\
    \Omega^b_{s,s}
        &=\frac{\rho^b_s\rho^b_s}{\rho^b_e}=4.
    \label{eq:fbc25_positive_block_index_values}
\end{align}
Equations~\ref{eq:fbc25_positive_block_index_identity_axes},
\ref{eq:fbc25_positive_block_index_value_a} and
\ref{eq:fbc25_positive_block_index_values} exhaust the four pairs $(g,h)$ in
each block, so $\Omega$ is positive.

The cocycle condition $d\Omega=d(d\rho)=1$ holds for the generalized
coboundary of any one-cochain, with no hypothesis on the action: substituting
Eq.~\ref{eq:fbc25_positive_block_index_coboundary} into both sides of the
cocycle equation gives
\begin{align}
    (d\rho)^x_{g,hk}(d\rho)^x_{h,k}
        &=\frac{\rho^{(hk)x}_g\rho^x_{hk}}{\rho^x_{ghk}}
          \cdot\frac{\rho^{k x}_h\rho^x_k}{\rho^x_{hk}}
        =\frac{\rho^{(hk)x}_g\rho^{k x}_h\rho^x_k}{\rho^x_{ghk}},
    \label{eq:fbc25_positive_block_index_cocycle_left}\\
    (d\rho)^{k x}_{g,h}(d\rho)^x_{gh,k}
        &=\frac{\rho^{h(k x)}_g\rho^{k x}_h}{\rho^{k x}_{gh}}
          \cdot\frac{\rho^{k x}_{gh}\rho^x_k}{\rho^x_{ghk}}
        =\frac{\rho^{h(k x)}_g\rho^{k x}_h\rho^x_k}{\rho^x_{ghk}},
    \label{eq:fbc25_positive_block_index_cocycle_right}
\end{align}
and the action law $h(k x)=(hk)x$ makes
Eq.~\ref{eq:fbc25_positive_block_index_cocycle_left} and
Eq.~\ref{eq:fbc25_positive_block_index_cocycle_right} equal.\footnote{Formal
declaration:
\leanid{TNLean.Algebra.ActionTensorGauge.isGeneralizedCocycle\_coboundary} in
\doclink{TNLean/Algebra/GeneralizedCocycle}.}

Thus $\Omega$ is a positive generalized $2$-cocycle and a generalized
coboundary, and it satisfies the hypotheses of the corollary. If a primitive
$\chi_g$ were independent of the block, then $(d\chi)^x_{s,s}$ would also be
independent of $x$, contradicting
Eq.~\ref{eq:fbc25_positive_block_index_values}.

\section{Corrected statement}

The primitive in the corollary must retain its block label:
\begin{align}
    \Omega^x_{g,h}
        &=(d\chi)^x_{g,h}
          =\frac{\chi^{h x}_g\chi^x_h}{\chi^x_{gh}}.
    \label{eq:fbc25_positive_block_index_corrected}
\end{align}
This agrees with the determinant construction immediately preceding the
corollary and leaves all other hypotheses and conclusions unchanged.

\section{Formal treatment}

The formal theorem quantifies over a block-dependent cochain and proves
Eq.~\ref{eq:fbc25_positive_block_index_corrected}.\footnote{Formal
declarations: \leanid{TNLean.Algebra.ActionTensorGauge.coboundary} and
\leanid{TNLean.Algebra.LSymbol.exists\_positive\_coboundary} in
\doclink{TNLean/Algebra/PositiveGeneralizedCocycle}.}

\textbf{Verdict: resolved local correction.}
The surrounding determinant construction and the explicit two-block example
show that the block superscript on $\chi$ is necessary.

\bibliographystyle{alpha}
\bibliography{references}

\end{document}